\input{preamble.tex}
\title{The semidirect construction as a join}
\author{Human Compression Project}\date{Private draft, 11 October 2026}
\begin{document}\maketitle
Let $\alpha=(A,(A_i)_{i\in S})$ and
$\beta=(B,(B_j)_{j\in T})$ be flag-transitive coset geometries of
finite rank. Write $H_A=A_S$ and $H_B=B_T$. Suppose that $B$ acts on
$A$ by automorphisms and that an indexed action on $S$ is specified,
with $b(A_i)=A_{bi}$. This specification also covers coincident
parabolics: an action on an unlabelled set of subgroups alone does not
automatically specify an action on their repeated indices. Let
$\mathcal O$ be the orbit partition of $S$.

For an orbit $s\in\mathcal O$, put $A_s=\bigcap_{i\in s}A_i$.
In $G=A\rtimes B$ take the indexed parabolics
\[
 P_s=A_s\rtimes B\quad(s\in\mathcal O),\qquad
 P_j=A\rtimes B_j\quad(j\in T).
\]
This is the first semidirect construction in Section 4.1 of
Piedade--Tranchida, arXiv:2505.24662v1. The observation below explains
its incidence without a separate product-set computation.

\begin{proposition}\label{prop:join}
Let $\overline\alpha=(A,(A_s)_{s\in\mathcal O})$. The coset
incidence system $(G,(P_i))$ is the join
$\overline\alpha*\beta$: the two component incidences are unchanged
and every element of one component is incident with every element of
the other. Both $\overline\alpha$ and the join are flag-transitive
geometries. For $J=J_O\sqcup J_T$, put
$E=\bigcup_{s\in J_O}s$. Then
\[
 G_J=A_E\rtimes B_{J_T},\qquad G_{\mathcal O\sqcup T}=H_A\rtimes H_B.
\]
\end{proposition}
\begin{proof}
Every $A_s$ is $B$-invariant. In semidirect coordinates,
\[
 (a,b)P_s=(aA_s)\times B,\qquad
 (a,b)P_j=A\times(bB_j).
\]
Here the second equality uses that $b$ acts bijectively on $A$.
Thus an orbit-type element is determined by the coset $aA_s$, a
$T$-type element by $bB_j$, and their incidences are exactly those
stated. The parabolic formula follows by intersecting the two
coordinates; this is ordinary intersection of Cartesian products.

The cosets of the $A_s$ belong to the family of cosets of standard
parabolics of $\alpha$. That family is Helly by flag transitivity:
expand a parabolic coset into the maximal-parabolic cosets whose
intersection defines it and apply flag transitivity to the expanded
flag. Hence $A$ is flag-transitive on $\overline\alpha$, and every
flag there extends to a chamber. The same is assumed for $\beta$.
For a flag of the join, choose $b\in B$ carrying its $\beta$ part
to the standard flag. This acts on its $\overline\alpha$ part;
then an element of $A$ carries that part to the standard flag while
fixing the $\beta$ part. Thus $G$ is flag-transitive. Alternatively,
a common $B$ coordinate and a common $A$ coordinate in the displayed
cosets give a common group representative directly.
\end{proof}

\begin{proposition}\label{prop:residues}
If both inputs are residually connected, so is the semidirect join.
Its panel indices are
\[
 [A_{S\setminus s}:H_A]\quad(s\in\mathcal O),\qquad
 [B_{T\setminus\{j\}}:H_B]\quad(j\in T).
\]
It is firm when both inputs are firm. If both inputs are thin, the
join is thin exactly when every orbit in $S$ is a singleton.
\end{proposition}
\begin{proof}
For clarity we first prove residual connectivity of
$\overline\alpha$. Put $R_i=A_{S\setminus\{i\}}$.
Residual connectivity of the finite-rank flag-transitive input gives
\[
 A_{S\setminus E}=\langle H_A,R_i:i\in E\rangle
 \qquad(E\subseteq S).
\]
For $|E|=0$ or $1$ this is the definition. In higher rank it follows
inductively from the residue-generation criterion. Take a coarse
residue with at least two orbit types remaining, whose union is $E$.
For every $i\in E$ choose a remaining orbit $s$ not containing $i$.
Then $R_i\leq A_{(S\setminus E)\cup s}$, one of the maximal
parabolics of that residue. Those parabolics also contain $H_A$ and
therefore generate its whole group $A_{S\setminus E}$. This proves
residual connectivity of the coarse geometry.

A residue of a join is the join of the corresponding component
residues. If both components have a type remaining, the incidence
graph is connected because all vertices across the two components
are incident. If just one component has types remaining, residual
connectivity is that of the corresponding input residue. This proves
the first assertion, including empty component type sets.

The panel formulas follow by omitting just one type from the
parabolic formula in Proposition~\ref{prop:join}. For an orbit $s$,
choose $i\in s$. Firmness of $\alpha$ gives $H_A<R_i$, and
$R_i\leq A_{S\setminus s}$, so its panel index is at least two.
The other panels have exactly their $\beta$ sizes.

Suppose both inputs are thin. Singleton orbit indices are two. If
$s$ contains distinct $i,k$, choose $x\in R_i\setminus H_A$ and
$y\in R_k\setminus H_A$. We have $R_i\cap R_k=H_A$ by their
definitions. The four left cosets
\[
 H_A,\quad xH_A,\quad yH_A,\quad xyH_A
\]
are distinct. For example, $xyH_A=xH_A$ would put $y$ in $H_A$;
$xyH_A=yH_A$ would put $x\in yH_Ay^{-1}\leq R_k$, a
contradiction; and $xyH_A=H_A$ would put $x\in R_k$.
Thus $[A_{S\setminus s}:H_A]\geq4$, ruling out thinness.
The argument uses only subgroup intersections and does not assume
that $H_A$ is normal in all of $A$.
\end{proof}

The source's statement that this construction is ``never thin'' is
therefore too strong. With trivial type action, it is the join of
the original thin geometries and is thin. What obstructs thinness
is an orbit containing at least two old types. In particular, the
semidirect construction for an infinite cyclic stable letter in the
HNN example has an infinite new panel because its rank-one cyclic
input already has one; that example does not show that all finite
semidirect joins fail thinness.

The common twisting theorem in the companion orbit paper handles the
second construction of source Section 4.2. Its exact Borel-dependent
panel formula explains the different thinness condition there. The
join proof above retains the source construction and replaces its
coordinate calculations by a direct incidence identification. No
historical priority, full-source recovery or numerical factor is
claimed by this private note.
\end{document}
